\documentclass[11pt]{article}
\usepackage{amsmath,amssymb,amsthm}
\usepackage[margin=2.8cm]{geometry}

\newtheorem{theorem}{Theorem}
\newtheorem{proposition}[theorem]{Proposition}
\newtheorem{lemma}[theorem]{Lemma}
\theoremstyle{definition}
\newtheorem{definition}[theorem]{Definition}
\theoremstyle{remark}
\newtheorem{remark}[theorem]{Remark}

\newcommand{\abs}[1]{\lvert #1\rvert}
\newcommand{\norm}[2]{\lVert #1\rVert_{#2}}
\newcommand{\cJ}{\mathcal J}
\newcommand{\Z}{\mathbb Z}
\newcommand{\W}{W_N}

\title{The finite \(J\)-norm is dominated by the discrete \(K\)-norm\\
\large (the direction \(\cJ \lesssim \norm{\cdot}{K}\) of Janson's Lemma~1, over a monoid)}
\date{}

\begin{document}
\maketitle

\section{Setting and statement}

Let \(\alpha\) be an additive commutative monoid with a preorder \(\le\) and a map
\(x \mapsto \abs{x}\) from \(\alpha\) to itself satisfying
\begin{equation}\label{eq:abs-add}
  \abs{x+y} \le \abs{x} + \abs{y} \qquad (x, y \in \alpha),
\end{equation}
together with the following \emph{refined Riesz decomposition}:
\begin{itemize}
  \item[(R)] if \(\abs{x} \le \abs{y} + \abs{z}\), then there are \(x_0, x_1 \in \alpha\) with
  \[
    x = x_0 + x_1,\qquad \abs{x_0} \le \abs{y},\quad \abs{x_1} \le \abs{z},\qquad
    \abs{x_0} \le \abs{x},\quad \abs{x_1} \le \abs{x}.
  \]
\end{itemize}
The first three conditions of (R) are the field \texttt{exists\_decomp} of
\texttt{Abs.IsModulus}; the last two are new. They hold for the three existing instances with
the witnesses already used there: on \(\mathbb R\), \(x_0\) is \(x\) clamped to
\([-\abs{y}, \abs{y}]\) and \(x_1 = x - x_0\) has the sign of \(x\); on \([0,\infty]\),
\(x_0 = \min(x, y)\) and \(x_1 = x - x_0\); on products, (R) is checked pointwise.

Let \(A_0, A_1\) be quasinorms on \(\alpha\) with values in \([0,\infty]\) that are
\emph{solid}: \(\abs{x} \le \abs{y}\) implies \(\norm{x}{i} \le \norm{y}{i}\), where
\(\norm{\cdot}{i}\) is the quasinorm of \(A_i\). No other property of \(A_0, A_1\) is used; in
particular their quasi-triangle constants never appear.

Fix \(0 < \theta < 1\) and \(0 < q \le \infty\). For \(\nu \in \Z\), \(x, u \in \alpha\) put
\[
  K_\nu(x) = \inf_{x = x_0 + x_1} \bigl(\norm{x_0}{0} + 2^\nu \norm{x_1}{1}\bigr),
  \qquad
  J_\nu(u) = \max\bigl(\norm{u}{0},\ 2^\nu \norm{u}{1}\bigr),
\]
and \(\Delta = \{x \in \alpha : \norm{x}{0} < \infty,\ \norm{x}{1} < \infty\}\).
For a family \(f = (f_\nu)_{\nu \in S}\) in \([0,\infty]\) indexed by a countable set \(S\),
\(\norm{f}{\ell^q}\) is \((\sum_{\nu\in S} f_\nu^q)^{1/q}\) if \(q < \infty\) and
\(\sup_{\nu \in S} f_\nu\) if \(q = \infty\). The two norms to be compared are
\[
  \norm{x}{K} = \bigl\lVert (2^{-\nu\theta} K_\nu(x))_{\nu \in \Z} \bigr\rVert_{\ell^q}
  \qquad (\texttt{discreteKMethod}),
\]
\[
  \cJ(x) = \inf \Bigl\{ \bigl\lVert (2^{-\nu\theta} J_\nu(u_\nu))_{\nu \in \W} \bigr\rVert_{\ell^q}
  : N \in \mathbb N,\ u \colon \W \to \alpha,\ \textstyle\sum_{\nu \in \W} u_\nu = x \Bigr\}
  \qquad (\texttt{jInfNorm}),
\]
where \(\W = \{-N, \dots, N-1\}\). In Lean the window is indexed by \(k \in
\texttt{Fin}\,(2N)\) with \(\nu = k - N\).

\begin{theorem}\label{thm:main}
  For every \(x \in \Delta\),
  \[
    \cJ(x) \le 2^{\min(\theta,\, 1-\theta)}\, \norm{x}{K} .
  \]
\end{theorem}

The constant is at most \(\sqrt 2\) and does not depend on \(q\) or on the couple.

\paragraph{Idea of the proof.}
For each level \(\mu\) choose a near-optimal decomposition \(x = y_\mu + z_\mu\) for
\(K_\mu(x)\). Using (R), split \(x\) into \(2N\) pieces, the piece \(u_\nu\) being dominated
by \(y_{\nu+1}\) and by \(z_\nu\). Then \(J_\nu(u_\nu) \lesssim K_{\nu+1}(x)\), except for one
of the two norms of each of the two end pieces, which is controlled only by \(x\) itself.
After weighting, these end terms decay geometrically in \(N\), because \(0<\theta<1\) and
\(x \in \Delta\). The \(\ell^q\)-norm of a maximum is estimated by a sum of \(q\)-th powers,
so no quasi-triangle inequality in \(\ell^q\) is needed.

\section{Three elementary facts}

\begin{lemma}\label{lem:K-dyadic}
  For every \(x \in \alpha\) and \(\nu \in \Z\),
  \(K_\nu(x) \le K_{\nu+1}(x) \le 2 K_\nu(x)\).
\end{lemma}

\begin{proof}
  For any decomposition \(x = x_0 + x_1\),
  \[
    \norm{x_0}{0} + 2^\nu\norm{x_1}{1}
    \le \norm{x_0}{0} + 2^{\nu+1}\norm{x_1}{1}
    \le 2\bigl(\norm{x_0}{0} + 2^\nu\norm{x_1}{1}\bigr).
  \]
  Taking the infimum over all decompositions gives both inequalities.
\end{proof}

No finiteness assumption on \(x\) is needed; the hypothesis of \texttt{monotone\_kNorm} in
\texttt{Basic.lean} is superfluous for dyadic points. The second inequality is a case of
\texttt{kNorm\_le\_mul}. The main proof uses only the first one.

\begin{lemma}\label{lem:near-opt}
  For every \(x \in \alpha\), \(\mu \in \Z\) and \(\delta > 0\) there is a decomposition
  \(x = y + z\) with \(\norm{y}{0} + 2^\mu \norm{z}{1} \le K_\mu(x) + \delta\).
\end{lemma}

\begin{proof}
  If \(K_\mu(x) = \infty\), any decomposition works, e.g.\ \(x = 0 + x\). Otherwise
  \(K_\mu(x) < K_\mu(x) + \delta\), and by the definition of the infimum some decomposition
  has cost below \(K_\mu(x) + \delta\) (\texttt{exists\_decomp\_lt\_of\_lt\_kNorm}).
\end{proof}

An additive error is used because the infimum need not be attained, and a multiplicative error
\(1 + \varepsilon\) is useless when \(K_\mu(x) = 0\).

\begin{lemma}\label{lem:lq-max}
  Let \(f, g\) be families in \([0,\infty]\) indexed by a countable set \(S\). If
  \(q < \infty\), then
  \(\norm{\max(f,g)}{\ell^q}^q \le \norm{f}{\ell^q}^q + \norm{g}{\ell^q}^q\).
  If \(q = \infty\), then \(\norm{\max(f,g)}{\ell^\infty} = \max(\norm{f}{\ell^\infty},
  \norm{g}{\ell^\infty})\).
\end{lemma}

\begin{proof}
  For \(q < \infty\), sum the pointwise inequality \(\max(s,t)^q \le s^q + t^q\) over \(S\).
  For \(q = \infty\), both sides are the supremum of all values of \(f\) and \(g\).
\end{proof}

\section{Splitting \(x\) into \(2N\) pieces}

This is the only place where (R) is used. The pieces are split off one at a time from the
part of \(x\) not yet assigned, so no subtraction is needed.

\begin{proposition}\label{prop:split}
  Let \(N \ge 1\), \(x \in \alpha\), and for \(\mu = -N+1, \dots, N-1\) let
  \(x = y_\mu + z_\mu\) be arbitrary decompositions. Then there are
  \(u_\nu \in \alpha\), \(\nu \in \W\), with \(\sum_{\nu \in \W} u_\nu = x\) and, for every
  \(\nu \in \W\),
  \begin{enumerate}
    \item[(a)] \(\abs{u_\nu} \le \abs{y_{\nu+1}}\) if \(\nu \le N-2\);
    \item[(b)] \(\abs{u_\nu} \le \abs{z_\nu}\) if \(\nu \ge -N+1\);
    \item[(c)] \(\abs{u_\nu} \le \abs{x}\).
  \end{enumerate}
\end{proposition}

\begin{proof}
  We construct remainders \(r_{-N}, \dots, r_{N-1}\) and pieces \(u_{-N}, \dots, u_{N-2}\)
  recursively, maintaining for each \(\nu\) the invariant
  \begin{equation}\label{eq:inv}
    x = \sum_{-N \le \kappa < \nu} u_\kappa + r_\nu, \qquad \abs{r_\nu} \le \abs{x}.
  \end{equation}
  Start with \(r_{-N} = x\); then \eqref{eq:inv} holds by reflexivity, the sum being empty.

  \emph{Recursive step.} Let \(-N \le \nu \le N-2\) and suppose \(r_\nu\) satisfies
  \eqref{eq:inv}. Since \(x = y_{\nu+1} + z_{\nu+1}\), inequality \eqref{eq:abs-add} and
  transitivity give
  \[
    \abs{r_\nu} \le \abs{x} \le \abs{y_{\nu+1}} + \abs{z_{\nu+1}} .
  \]
  Apply (R) to \(r_\nu\). It gives \(r_\nu = u_\nu + r_{\nu+1}\) with
  \begin{equation}\label{eq:step}
    \abs{u_\nu} \le \abs{y_{\nu+1}}, \qquad \abs{r_{\nu+1}} \le \abs{z_{\nu+1}}, \qquad
    \abs{u_\nu} \le \abs{r_\nu}, \qquad \abs{r_{\nu+1}} \le \abs{r_\nu}.
  \end{equation}
  Substituting \(r_\nu = u_\nu + r_{\nu+1}\) into \eqref{eq:inv} and using associativity gives
  the identity in \eqref{eq:inv} for \(\nu+1\). The last inequality in \eqref{eq:step} and
  \(\abs{r_\nu} \le \abs{x}\) give \(\abs{r_{\nu+1}} \le \abs{x}\).

  \emph{Last piece.} Put \(u_{N-1} = r_{N-1}\). The invariant \eqref{eq:inv} for
  \(\nu = N-1\) reads \(x = \sum_{\nu \in \W} u_\nu\).

  \emph{The bounds.} Property (a) is the first inequality of \eqref{eq:step}. For (b) and (c)
  note that \(\abs{u_\nu} \le \abs{r_\nu}\) for all \(\nu \in \W\): for \(\nu \le N-2\) by
  \eqref{eq:step}, and for \(\nu = N-1\) by definition. Now \(\abs{r_\nu} \le \abs{x}\) by
  \eqref{eq:inv}, which gives (c); and for \(\nu \ge -N+1\), \(\abs{r_\nu} \le \abs{z_\nu}\) by
  the second inequality of \eqref{eq:step} at step \(\nu - 1\), which gives (b).
\end{proof}

Both new conditions of (R) are needed. The condition \(\abs{x_1} \le \abs{x}\) keeps the
remainders below \(\abs{x}\), without which (R) could not be applied at the next step; the
condition \(\abs{x_0} \le \abs{x}\) passes the domination of the remainder by \(z_\nu\) to the
new piece.

\section{The pointwise estimate}

\begin{lemma}\label{lem:pointwise}
  Let \(x \in \alpha\), \(\delta > 0\) and \(N \ge 1\). For \(\mu = -N+1, \dots, N-1\) choose
  decompositions \(x = y_\mu + z_\mu\) as in Lemma~\ref{lem:near-opt}, i.e.
  \begin{equation}\label{eq:near-opt}
    \norm{y_\mu}{0} + 2^\mu \norm{z_\mu}{1} \le K_\mu(x) + \delta,
  \end{equation}
  and let \((u_\nu)_{\nu \in \W}\) be given by Proposition~\ref{prop:split}. Put
  \(E_{-N} = 2^{-N}\norm{x}{1}\), \(E_{N-1} = \norm{x}{0}\) and \(E_\nu = 0\) for the other
  \(\nu \in \W\). Then for every \(\nu \in \W\),
  \[
    J_\nu(u_\nu) \le \max\bigl(K_{\nu+1}(x) + \delta,\ E_\nu\bigr).
  \]
\end{lemma}

\begin{proof}
  Since \(N \ge 1\), the indices \(-N\) and \(N-1\) are distinct, so each \(\nu\) is
  exceptional in at most one of the two entries of \(J_\nu(u_\nu)\). We bound both entries
  by the right-hand side; solidity turns each modulus inequality of
  Proposition~\ref{prop:split} into an inequality of quasinorms.

  \emph{The entry \(\norm{u_\nu}{0}\).} If \(\nu \le N-2\), then by (a) and
  \eqref{eq:near-opt} at \(\mu = \nu+1\),
  \(\norm{u_\nu}{0} \le \norm{y_{\nu+1}}{0} \le K_{\nu+1}(x) + \delta\).
  If \(\nu = N-1\), then by (c), \(\norm{u_\nu}{0} \le \norm{x}{0} = E_{N-1}\).

  \emph{The entry \(2^\nu \norm{u_\nu}{1}\).} If \(\nu \ge -N+1\), then by (b),
  \eqref{eq:near-opt} at \(\mu = \nu\), and Lemma~\ref{lem:K-dyadic},
  \(2^\nu\norm{u_\nu}{1} \le 2^\nu\norm{z_\nu}{1} \le K_\nu(x) + \delta \le K_{\nu+1}(x) +
  \delta\). If \(\nu = -N\), then by (c), \(2^{-N}\norm{u_\nu}{1} \le 2^{-N}\norm{x}{1} =
  E_{-N}\).
\end{proof}

\section{Proof of Theorem~\ref{thm:main}}

We prove \(\cJ(x) \le 2^\theta \norm{x}{K}\); the bound with \(2^{1-\theta}\) is
Remark~\ref{rem:other-constant}. Fix \(x \in \Delta\), \(N \ge 1\) and \(\delta > 0\), and let
\(u\) be as in Lemma~\ref{lem:pointwise}. Write
\[
  a_\mu = 2^{-\mu\theta} K_\mu(x), \qquad
  b_\nu = 2^{-\nu\theta} J_\nu(u_\nu), \qquad
  e_\nu = 2^{-\nu\theta} E_\nu .
\]
Since \(2^{-\nu\theta} = 2^\theta \cdot 2^{-(\nu+1)\theta}\), multiplying
Lemma~\ref{lem:pointwise} by \(2^{-\nu\theta}\) gives
\begin{equation}\label{eq:weighted}
  b_\nu \le \max\bigl(2^\theta (a_{\nu+1} + \delta\, 2^{-(\nu+1)\theta}),\ e_\nu\bigr)
  \qquad (\nu \in \W).
\end{equation}
The end terms are
\[
  e_{-N} = 2^{N\theta} 2^{-N}\norm{x}{1} = 2^{-(1-\theta)N}\norm{x}{1},
  \qquad
  e_{N-1} = 2^{-(N-1)\theta}\norm{x}{0}.
\]
Since \(\sum_\nu u_\nu = x\), the family \(u\) is admissible in the definition of \(\cJ\), so
\(\cJ(x) \le \norm{b}{\ell^q}\).

\emph{Case \(q < \infty\).} By \eqref{eq:weighted}, Lemma~\ref{lem:lq-max}, and the fact that
\(e\) vanishes except at \(-N\) and \(N-1\),
\[
  \cJ(x)^q \le \sum_{\nu \in \W} b_\nu^q
  \le 2^{\theta q} \sum_{\nu \in \W} \bigl(a_{\nu+1} + \delta\, 2^{-(\nu+1)\theta}\bigr)^q
    + e_{-N}^q + e_{N-1}^q .
\]
The left-hand side does not depend on \(\delta\). The right-hand side is a finite sum of
functions of \(\delta\) that are continuous on \([0,\infty)\) with values in \([0,\infty]\)
(addition, multiplication by a finite constant, and \(t \mapsto t^q\) are continuous there).
Letting \(\delta \to 0\) and substituting \(\mu = \nu + 1\),
\[
  \cJ(x)^q \le 2^{\theta q} \sum_{\mu = -N+1}^{N} a_\mu^q + e_{-N}^q + e_{N-1}^q
  \le 2^{\theta q}\norm{x}{K}^q + e_{-N}^q + e_{N-1}^q ,
\]
where the last step drops the terms with \(\mu \notin \{-N+1, \dots, N\}\). Since
\(0 < \theta < 1\) and \(\norm{x}{0}, \norm{x}{1} < \infty\), both \(e_{-N}\) and
\(e_{N-1}\) tend to \(0\) as \(N \to \infty\), and so do their \(q\)-th powers. Letting
\(N \to \infty\) gives \(\cJ(x)^q \le 2^{\theta q}\norm{x}{K}^q\), and taking \(q\)-th roots
proves the claim.

\emph{Case \(q = \infty\).} By \eqref{eq:weighted} and Lemma~\ref{lem:lq-max}, and since
\(2^{-(\nu+1)\theta} \le 2^{(N-1)\theta}\) on \(\W\),
\[
  \cJ(x) \le \max_{\nu \in \W} b_\nu
  \le \max\bigl(2^\theta(\norm{x}{K} + \delta\, 2^{(N-1)\theta}),\ e_{-N},\ e_{N-1}\bigr).
\]
Let \(\delta \to 0\), then \(N \to \infty\). \qed

\medskip
The two limits are taken one after the other: \(\delta \to 0\) for fixed \(N\), then
\(N \to \infty\). This is legitimate because \(\cJ(x)\) depends on neither, and the end terms
do not depend on \(\delta\). No assumption \(\norm{x}{K} < \infty\) is needed.

\section{Remarks}

\begin{remark}[The constant \(2^{1-\theta}\)]\label{rem:other-constant}
  Feed Proposition~\ref{prop:split} the decompositions one level lower: for
  \(\mu = -N+1, \dots, N-1\), let \(x = y_\mu + z_\mu\) satisfy \eqref{eq:near-opt} at level
  \(\mu - 1\). For \(\nu \le N-2\), (a) and Lemma~\ref{lem:K-dyadic} give
  \(\norm{u_\nu}{0} \le K_\nu(x) + \delta \le 2(K_{\nu-1}(x) + \delta)\); for
  \(\nu \ge -N+1\), (b) gives
  \(2^\nu\norm{u_\nu}{1} = 2 \cdot 2^{\nu-1}\norm{u_\nu}{1} \le 2(K_{\nu-1}(x) + \delta)\).
  The exceptional entries are unchanged, so
  \(J_\nu(u_\nu) \le \max(2(K_{\nu-1}(x) + \delta), E_\nu)\). Since
  \(2 \cdot 2^{-\nu\theta} = 2^{1-\theta} 2^{-(\nu-1)\theta}\), the rest of the proof goes
  through with \(2^\theta a_{\nu+1}\) replaced by \(2^{1-\theta} a_{\nu-1}\). This uses the
  second inequality of Lemma~\ref{lem:K-dyadic}. Choosing the better of the two versions gives
  the constant \(2^{\min(\theta, 1-\theta)}\).
\end{remark}

\begin{remark}[Where the hypotheses are used]
  \leavevmode
  \begin{itemize}
    \item (R), \eqref{eq:abs-add} and transitivity of \(\le\): only in
      Proposition~\ref{prop:split}.
    \item Solidity of \(A_0, A_1\): only in Lemma~\ref{lem:pointwise}.
    \item \(x \in \Delta\) and \(0 < \theta < 1\): only in \(e_{-N}, e_{N-1} \to 0\). Each end
      uses one half of each condition: \(e_{-N}\) uses \(\theta < 1\) and
      \(\norm{x}{1} < \infty\), \(e_{N-1}\) uses \(\theta > 0\) and \(\norm{x}{0} < \infty\).
    \item The quasi-triangle inequalities of \(A_0\), \(A_1\) and of \(\ell^q\): never.
      Lemma~\ref{lem:lq-max} replaces Minkowski's inequality, which is why the constant does
      not depend on \(q\).
  \end{itemize}
  For \(q = 0\), Mathlib's \texttt{eLpNorm} is identically \(0\), so \(\cJ \equiv 0\) and the
  inequality is trivial.
\end{remark}

\begin{remark}[Why \(x \in \Delta\) is necessary]
  Let \(q \neq 0\). If \(\cJ(x) < \infty\), some admissible \(\ell^q\)-norm is finite, so each
  \(J_\nu(u_\nu)\) is finite and each \(u_\nu\) lies in \(\Delta\). Since \(\Delta\) is an
  additive submonoid (by the quasi-triangle inequality with finite constants), the finite sum
  \(x\) lies in \(\Delta\). Hence \(\cJ = \infty\) outside \(\Delta\), whereas
  \(\norm{x}{K}\) may be finite there.
\end{remark}

\end{document}
